\documentclass{article}
\usepackage[a4paper, top=2cm, left=2cm, right=2cm, bottom=1cm]{geometry}
\usepackage{fontspec}
\usepackage{tikz-cookingsymbols}
\usepackage{tabularx}
\usepackage{multirow}
\usepackage{enumitem}
\usepackage{array}

\setlist[itemize]{leftmargin=*,nosep}
\renewcommand{\familydefault}{\sfdefault}
\setlength{\parindent}{0pt}
\setlength{\parskip}{0.5em}

\newlength{\colA}
\newlength{\colB}
\setlength{\colA}{0.07\textwidth}
\setlength{\colB}{0.15\textwidth}
\newcolumntype{A}{>{\raggedright\arraybackslash}p{\colA}}
\newcolumntype{B}{>{\raggedright\arraybackslash}p{\colB}}
\newcolumntype{C}{>{\raggedright\arraybackslash}X}

\begin{document}
\pagestyle{empty}

{\Huge\bfseries Gâteau aux noix}\par
\vspace{2cm}
\noindent
\begin{minipage}[t]{0.3\textwidth}
\vspace{0pt}%
\rule{\linewidth}{1.5pt}
\TopBottomHeat\quad{} temps total \textbf{140}~min \\
\Grill\quad{} préparation \textbf{50}~min \\
\ConvectionOven\quad{} cuisson \textbf{40}~min \\
\rule{\linewidth}{1.5pt}

\subsection*{Description}
Ce gâteau aux noix séduit dès la première bouchée par son équilibre parfait
entre une pâte sablée délicatement friable et une farce généreuse mêlant
noix croquantes, caramel doré et miel parfumé. Doux sans être lourd,
riche sans être écœurant, il offre une texture fondante relevée par
le croustillant des noix et la profondeur du caramel. Servi tiède
ou à température ambiante, il dégage un parfum chaleureux et
gourmand qui évoque les desserts traditionnels d'automne, tout en
offrant une élégance rustique irrésistible. Une véritable tentation
pour tous les amateurs de saveurs authentiques et de pâtisseries raffinées.

\subsection*{Conseils}

Se conserve env. 2 semaines au réfrigérateur bien emballé.

\end{minipage}
\hfill
\begin{minipage}[t]{0.68\textwidth}
\vspace{0pt}%
\textbf{Préparation du moule}\vskip 0.75em

\begin{itemize}
  \item Préparer le moule.
\end{itemize}

\vspace{1em}
\textbf{Pâte sablée}\vskip 0.75em

\begin{tabularx}{\textwidth}{@{}A B C@{}}
    150~g &
    farine &
      Mettre dans un saladier.
\\
    100~g &
    sucre &
\\
    1~pincée &
    sel &
\\
\multicolumn{3}{@{}l}{\rule{\textwidth}{0.2pt}}\\
    100~g &
    beurre froid &
      Couper en morceaux, ajouter. \newline < Travailler soigneusement à la main jusqu'à ce que la préparation soit friable.
\\
\multicolumn{3}{@{}l}{\rule{\textwidth}{0.2pt}}\\
    1 &
    œuf &
      Ajouter. \newline Travailler rapidement à la main, ne pas pétrir. \newline Couvrir et mettre au réfrigérateur 30 minutes.
\\
\end{tabularx}

\vspace{1em}
\textbf{Farce}\vskip 0.75em

\begin{tabularx}{\textwidth}{@{}A B C@{}}
    130~g &
    noix &
      Hacher grossièrement, réserver.
\\
\multicolumn{3}{@{}l}{\rule{\textwidth}{0.2pt}}\\
    100~g &
    sucre &
      Verser dans une grande casserole. \newline Chauffer à feu très vif jusqu'à dissolution du sucre, revenir à feu vif. \newline Tourner légèrement la casserole. \newline Faire caraméliser sans remuer jusqu'à ce que la couleur devienne noisette. \newline Retirer la casserole du feu.
\\
    2~cs &
    eau &
\\
\multicolumn{3}{@{}l}{\rule{\textwidth}{0.2pt}}\\
    1.25~dl &
    crème entière &
      Ajouter, mélanger à la cuillère de cuisine. \newline < Réchauffer à petit feu en remuant jusqu'à dissolution complète du caramel, porter à ébullition env. 10 min.
\\
\multicolumn{3}{@{}l}{\rule{\textwidth}{0.2pt}}\\
    1.5~cs &
    miel liquide &
      Ajouter, mélanger, laisser refroidir.
\\
\end{tabularx}

\vspace{1em}
\textbf{Assemblage}\vskip 0.75em

\begin{tabularx}{\textwidth}{@{}A B C@{}}
    un peu
 &
    farine &
      Partager la pâte en 3 portions régulières. \newline < Étaler 1 portion en un cercle d'environ 4 mm d'épaisseur sur un peu de farine. \newline Poser dans le moule préparé. \newline Piquer avec une fourchette.
\\
\multicolumn{3}{@{}l}{\rule{\textwidth}{0.2pt}}\\
    un peu
 &
    farine &
      Former un rouleau avec la deuxième portion de pâte. \newline Appuyer sur les bords du moule jusqu'à obtenir un bord d'environ 4 cm de hauteur.
\\
\multicolumn{3}{@{}l}{\rule{\textwidth}{0.2pt}}\\
  & & Répartir la farce sur le fond de pâte. \newline Replier le bord qui dépasse sur la farce. \newline < Étaler la dernière portion en un cercle d'environ 4 mm et poser sur la farce. \newline Piquer avec une fourchette. \newline Cuire dans le bas du four, 200°C, 40-50 min.\\
\end{tabularx}

\vspace{1em}
\end{minipage}

\end{document}
